\chapter{Lines and Planes in Three Dimensions}

\section{Points and Vectors in Space}
Points in three dimensions are written $(x, y, z)$. A vector in space has three
components, $\vec{v} = \langle a, b, c\rangle$. A vector gives a direction
and a distance, and it can move a point to a new point.

\section{Lines in Space}\index{line!in space}
A line in space is determined by one point $P_0 = (x_0, y_0, z_0)$ on it and a
\newterm{direction vector} $\vec{d} = \langle a, b, c\rangle$ parallel to it.
Every point on the line can be reached from $P_0$ by moving some multiple $t$ of
$\vec{d}$. This gives the \index{vector equation!line}\newterm{vector equation}
of the line:
$$\vec{r}(t) = \langle x_0, y_0, z_0\rangle + t\langle a, b, c\rangle$$
Writing the components separately gives the \newterm{parametric equations}:
$$x = x_0 + at, \qquad y = y_0 + bt, \qquad z = z_0 + ct$$

For example, the line through $(1, 2, 3)$ with direction $\langle 2, -1, 4\rangle$ is
$x = 1 + 2t$, $y = 2 - t$, $z = 3 + 4t$. When $t = 1$, the point is $(3, 1, 7)$.

\section{Planes}\index{plane!in space}
A plane is determined by one point $P_0$ on it and a \newterm{normal vector}
$\vec{n} = \langle a, b, c\rangle$, which is perpendicular to every direction in the
plane. A point $P = (x, y, z)$ is on the plane exactly when the vector from
$P_0$ to $P$ is perpendicular to $\vec{n}$:
$$\vec{n} \cdot \langle x - x_0, y - y_0, z - z_0\rangle = 0$$
Expanding this gives the linear equation
$$ax + by + cz = d, \qquad \text{where } d = ax_0 + by_0 + cz_0$$

For the plane through $(1, 1, 1)$ with normal $\langle 2, -1, 2\rangle$, we get
$d = 2 - 1 + 2 = 3$, so the plane is $2x - y + 2z = 3$.

\section{Distance from a Point to a Plane}\index{distance!point to plane}
The distance from a point $P_0 = (x_0, y_0, z_0)$ to the plane $ax + by + cz = d$ is
$$D = \frac{|ax_0 + by_0 + cz_0 - d|}{\sqrt{a^2 + b^2 + c^2}}$$
The numerator measures how far $P_0$ is from satisfying the plane's equation,
and the denominator scales that by the length of the normal vector.

For the origin and the plane $2x - y + 2z = 3$: the numerator is $|0 - 3| = 3$, and
the denominator is $\sqrt{4 + 1 + 4} = 3$. The distance is $1$.

\section{Where Lines and Planes Meet}
A line and a plane usually meet at one point. To find it, put the parametric
equations of the line into the equation of the plane, and solve for $t$. Then
substitute that $t$ back into the line.

\begin{Exercise}[label=lp1, title=A point on a line]
Find the point on the line $x = 1 + 2t$, $y = 2 - t$, $z = 3 + 4t$ when $t = -1$.
\end{Exercise}
\begin{Answer}[ref=lp1]
$(1 - 2, 2 + 1, 3 - 4) = (-1, 3, -1)$.
\end{Answer}

\begin{Exercise}[label=lp2, title=Equation of a plane]
Find the equation of the plane through $(1, 0, 0)$, $(0, 1, 0)$, and $(0, 0, 1)$.
\end{Exercise}
\begin{Answer}[ref=lp2]
The normal $\langle 1, 1, 1\rangle$ works, because it is perpendicular to both
$\langle -1, 1, 0\rangle$ and $\langle -1, 0, 1\rangle$ (the vectors from $(1,0,0)$ to the other two points).
Using the point $(1, 0, 0)$, $d = 1$, so the plane is $x + y + z = 1$.
\end{Answer}

\begin{Exercise}[label=lp3, title=Distance to a plane]
Find the distance from the origin to the plane $x + y + z = 1$.
\end{Exercise}
\begin{Answer}[ref=lp3]
$D = \frac{|0 - 1|}{\sqrt{3}} = \frac{1}{\sqrt{3}} = \frac{\sqrt{3}}{3} \approx 0.577$.
\end{Answer}
